\par\Needspace{6\baselineskip}\medskip
\begin{center}\hypertarget{native-243-article173}{}\pdfbookmark[2]{173.}{article-173}173.\end{center}
It is especially important that each geocentric position of the celestial body on which we intend to base the orbit should come from several observations, if possible, combined so that any errors cancel one another as far as possible, rather than from an isolated observation. Observations only a few days apart --- or even 15 or 20 days apart, as circumstances warrant --- should therefore yield a single position for the computation rather than as many distinct positions. This position is, so to speak, a mean of them all and thus admits of far greater precision than the individual observations considered separately. This procedure rests on the following principles.

The geocentric places of the celestial body calculated from approximate elements should differ little from the true places, and the differences between them should change only very slowly, so that over a few days they may be regarded as practically constant, or at least their variations may be regarded as proportional to the times. If, therefore, the observations were free of all error, the differences between the places observed at times $t$, $t'$, $t''$, $t'''$ etc. and those computed from the elements, i.\,e. the differences between observed and computed longitudes and latitudes, or between observed and computed right ascensions and declinations, would be practically equal or at least would increase or decrease uniformly and very slowly. Let the observed right ascensions at those times be, e.\,g., $a$, $a'$, $a''$, $a'''$ etc., and the computed ones $a+\delta$, $a'+\delta'$, $a''+\delta''$, $a'''+\delta'''$ etc.; then the differences $\delta$, $\delta'$, $\delta''$, $\delta'''$ etc. will differ from the true deviations due to the elements only insofar as the observations themselves are erroneous: if those deviations can be regarded as constant for all these observations, the quantities $\delta$, $\delta'$, $\delta''$, $\delta'''$ etc. will give as many different determinations of the same magnitude. It will therefore be appropriate to take the arithmetic mean of those determinations as its correct value, insofar as there is no reason to prefer one \SourceLeaf{244}or another. If, however, the individual observations do not appear to merit the same degree of precision, let us suppose that their respective degrees of precision are to be reckoned proportional to the numbers $e$, $e'$, $e''$, $e'''$ etc., i.\,e. that errors inversely proportional to these numbers could have been made in the observations with equal ease; then, according to principles to be given below, the most probable mean value will no longer be the simple arithmetic mean, but
\[
=\frac{ee\delta+e'e'\delta'+e''e''\delta''+e'''e'''\delta'''+\text{etc.}}
{ee+e'e'+e''e''+e'''e'''+\text{etc.}}\mathpunct{\cdot}
\]
Taking this mean value $=\Delta$, we may take as the true right ascensions, respectively, $a+\delta-\Delta$, $a'+\delta'-\Delta$, $a''+\delta''-\Delta$, $a'''+\delta'''-\Delta$ etc., and it will then be immaterial which one we use in the computation. If, however, the observations are separated by too great an interval of time, or if sufficiently approximate elements of the orbit are not yet known, so that their deviations cannot be treated as constant for all the observations, it is easy to see that the only difference is that the mean deviation so found must be referred to a certain mean time, rather than assumed to be common to all the observations. This mean time must be derived from the individual times in the same way that $\Delta$ is derived from the individual deviations, and is therefore in general the time
$\frac{eet+e'e't'+e''e''t''+e'''e'''t'''+\text{etc.}}{ee+e'e'+e''e''+e'''e'''+\text{etc.}}\mathpunct{\cdot}$
If we wish to attain the utmost precision, we must therefore compute the geocentric place for that same time from the elements and then remove the mean error $\Delta$ from it, to obtain the most accurate position possible: in most cases, however, it will amply suffice to refer the mean error to the observation nearest the mean time. What we have said here about right ascensions applies equally to declinations, or, if preferred, to longitudes and latitudes: it will always be better, however, to compare the right ascensions and declinations computed from the elements directly with those observed; for this gives us not only a quicker computation, especially if we use the methods set out in arts.\ 53---60, but also the advantage that incomplete observations can be used. Moreover, if everything were referred to longitudes and latitudes, an observation made correctly in right ascension and poorly in declination (or the reverse) might be spoiled in both coordinates and so become entirely useless. --- The degree of precision to be assigned to the mean so found will, according to principles to be explained shortly, be $=\surd(ee+e'e'+e''e''+e'''e'''+\text{etc.})$, so that four or nine equally \SourceLeaf{245}accurate observations are required if the mean is to have twice or three times the precision, and so forth.
